%%%%%%%% EEML 2026 Extended Abstract — List-style version (10-seed) %%%%%%%%%%
\documentclass{article}
\usepackage{microtype}
\usepackage{graphicx}
\usepackage{subfigure}
\usepackage{booktabs}
\usepackage{hyperref}
\newcommand{\theHalgorithm}{\arabic{algorithm}}

\usepackage[accepted]{icml2025}

\usepackage{amsmath}
\usepackage{amssymb}
\usepackage{mathtools}
\usepackage{amsthm}
\usepackage[capitalize,noabbrev]{cleveref}

\theoremstyle{plain}
\newtheorem{proposition}{Proposition}[section]

\icmltitlerunning{Graph Construction from a Statistical Vocabulary}

\begin{document}

\twocolumn[
\icmltitle{Graph Construction from a Statistical Vocabulary:\\Structured Inductive Bias for Relational Learning on Time Series}

\icmlsetsymbol{equal}{*}

\begin{icmlauthorlist}
\icmlauthor{[Author Name]}{usyd}
\icmlauthor{[Supervisor Name]}{usyd}
\end{icmlauthorlist}

\icmlaffiliation{usyd}{The University of Sydney, Australia}

\vskip 0.3in
]

\printAffiliationsAndNotice{}

\begin{abstract}
Graph neural networks for multivariate time series require a graph that is rarely given. We parameterise graph construction over a vocabulary of $K \approx 125$ named pairwise statistics --- causal, spectral, linear, information-theoretic --- computed via \texttt{pyspi} \citep{cliff2023}. A learned weight vector $\mathbf{w}$, jointly optimised with a message-passing network, selects which statistical notion of coupling is task-relevant. We prove that symmetric construction operators cannot distinguish Markov-equivalent topologies. On synthetic VAR(1) classification, the full vocabulary achieves 100\% (std $< 1$\%), while the oracle-best single statistic reaches 96\% but with catastrophic seed failures (std 24\%) and symmetric baselines remain at or below 59\%. The learned $\mathbf{w}$ recovers spectral Granger causality as the dominant coupling mode, matching the generating process.
\end{abstract}

%-----------------------------------------------------------------------
\section{Introduction}
\label{sec:intro}
%-----------------------------------------------------------------------

Message passing neural networks (MPNNs) derive their power from propagating information along relational structure \citep{gilmer2017,battaglia2018}. In MTS settings, this structure is unobserved and must be constructed. Two regimes dominate current practice.

\textbf{Fixing the graph from a single statistic} --- Pearson correlation in functional connectivity \citep{bullmore2009}, spatial distance in traffic networks \citep{li2018dcrnn} --- is a strong commitment to one mathematical notion of dependence. If the task-relevant coupling is directed, nonlinear, or frequency-specific, this is an expressiveness limitation of the graph construction operator, not a data limitation. \citet{liu2025} recently benchmarked 239 pairwise statistics for functional connectivity mapping, finding ``substantial quantitative and qualitative variation across FC methods'' --- empirical confirmation that the choice of construction statistic is consequential yet understudied.

\textbf{Learning the graph from latent embeddings} --- Neural Relational Inference \citep{kipf2018}, Graph WaveNet \citep{wu2019}, MTGNN \citep{wu2020} --- recovers flexibility at the cost of semantics. The learned adjacency is opaque: an edge weight of 0.7 encodes no information about \emph{what form} of dependence it represents. Such graphs cannot be validated against domain knowledge or compared across tasks. They also place the full burden on the training signal to discover relational structure that statistical methodology could supply as prior knowledge.

\citet{velickovic2021} argued that aligning neural architectures with algorithmic structure yields favourable sample complexity; \citet{velickovic2022} demonstrated that injecting algorithmic priors improves sample efficiency over learning from scratch. We apply this principle to graph construction: rather than re-learning what ``connectivity'' means from raw data, we supply a pre-specified vocabulary of pairwise statistics with known mathematical properties and let end-to-end learning identify which elements are task-relevant.

\textbf{Our approach.} For each ordered pair $(i, j)$ in an MTS, we precompute a $K$-dimensional descriptor $\mathbf{E}_{ij}$ whose coordinates are named pairwise statistics --- transfer entropy, spectral Granger causality, phase locking value, covariance, and $K \approx 125$ others --- drawn from six complementary mathematical families using \texttt{pyspi} \citep{cliff2023}. Graph construction becomes:
%
\begin{equation}
A_{ij} = \mathrm{softplus}(b + \mathbf{w}^\top \mathbf{E}_{ij}), \quad \text{top-}d \text{ sparsification per node}
\end{equation}
%
where $\mathbf{w} \in \mathbb{R}^K$ is jointly optimised with the downstream MPNN. This is a deliberate linear probe of the SPI descriptor space \citep[cf.][]{alain2017}: the claim is that the vocabulary is sufficiently well-structured for a linear function to identify task-relevant edges. Group lasso regularisation \citep{yuan2006} over SPI families encourages family-level selection, producing an interpretable \emph{statistical signature} of the task's relational structure.

%-----------------------------------------------------------------------
\section{Expressiveness of Graph Construction Operators}
\label{sec:theory}
%-----------------------------------------------------------------------

We formalise the intuition that the choice of construction statistic constrains what an MPNN can learn, analogous to Weisfeiler--Lehman expressiveness results for message passing \citep{xu2019,morris2019}.

\begin{proposition}
Consider three 3-node VAR(1) motifs over nodes $A, B, C$: chain ($A \to B \to C$), fork ($A \leftarrow B \to C$), and collider ($A \to B \leftarrow C$). All three share skeleton $A$--$B$--$C$. Chain and fork are Markov equivalent \citep{verma1990}: they encode identical conditional independence relations ($A \perp C \mid B$, and no other). Any graph construction operator using only symmetric statistics --- correlation, mutual information, coherence, or any function $f$ satisfying $f(X_i, X_j) = f(X_j, X_i)$ --- assigns identical edge weights to chain and fork, regardless of sample size or estimation quality. No MPNN operating on such a graph can exceed $2/3$ accuracy on the three-class task.
\end{proposition}

\emph{Proof sketch.} Two DAGs are Markov equivalent iff they share the same skeleton and v-structures \citep{verma1990}. Chain and fork share skeleton $\{A$--$B,\; B$--$C\}$ and have no v-structures. Any symmetric statistic $f(X_i, X_j)$ depends only on the joint distribution $P(X_i, X_j)$, which is identical for Markov-equivalent structures. The collider ($A \to B \leftarrow C$) has a v-structure and is distinguishable by all methods --- serving as an internal control. $\square$

Directed statistics break the chain--fork equivalence. In the chain, $\mathrm{SGC}(A \to B) \gg \mathrm{SGC}(B \to A)$; in the fork, the pattern differs. The asymmetric descriptor tensor $\mathbf{E}_{ij} \neq \mathbf{E}_{ji}$ provides the information needed for discrimination.

%-----------------------------------------------------------------------
\section{Method}
\label{sec:method}
%-----------------------------------------------------------------------

\textbf{SPI descriptor tensor.} Given $X \in \mathbb{R}^{M \times T}$ (z-scored), we compute $K$ pairwise statistics per ordered pair via \texttt{pyspi} \citep{cliff2023}, yielding $\mathbf{E} \in \mathbb{R}^{M \times M \times K}$. Statistics are partitioned into families: \textbf{causal} (transfer entropy \citep{schreiber2000}, Granger causality \citep{granger1969}, spectral GC \citep{geweke1982}, phase slope index), \textbf{spectral} (coherence, PLV, PLI), \textbf{linear} (covariance, precision, cross-correlation), \textbf{information} (mutual information), \textbf{distance} (DTW, Euclidean), and \textbf{rank} (Spearman, Kendall). Crucially, causal measures are asymmetric: $\mathbf{E}_{ij} \neq \mathbf{E}_{ji}$, enabling directed graph construction.

\textbf{Learned graph construction.} $A_{ij} = \mathrm{softplus}(b + \mathbf{w}^\top \mathbf{E}_{ij})$ with $K + 1$ learnable parameters. Sparsified to top-$d$ outgoing edges per node. Group lasso: $\mathcal{L} = \mathcal{L}_{\mathrm{task}} + \lambda_1 \|\mathbf{w}\|_1 + \lambda_g \sum_{g \in \mathcal{G}} \|\mathbf{w}_g\|_2$. Training uses Adam with LR warmup over 60 epochs and restart selection (best of 2 initialisations).

\textbf{Edge-attributed MPNN.} Retained edges carry $\mathbf{E}_{ij}$ as attributes. An edge network $\phi(\mathbf{E}_{ij}) = \mathrm{MLP}_\phi(\mathbf{E}_{ij})$ conditions messages:
%
\begin{equation}
\mathbf{m}_{ij} = \mathrm{MLP}_m([\mathbf{h}_j \| \mathbf{h}_i \| \phi(\mathbf{E}_{ij})]), \quad \mathbf{h}_i' = \mathbf{h}_i + \mathrm{LN}\!\left(\textstyle\sum_j A_{ij} \cdot \mathbf{m}_{ij}\right)
\end{equation}
%
The SPI vocabulary informs both topology (via $\mathbf{w}$) and message content (via $\phi$). Global mean+max pooling and a classifier produce graph-level predictions.

%-----------------------------------------------------------------------
\section{Experiments}
\label{sec:experiments}
%-----------------------------------------------------------------------

\textbf{Synthetic topology classification.} We generate $M{=}10$ node VAR(1) processes ($T{=}500$) with a 3-node directed motif ($A{\to}B{\to}C$, $A{\leftarrow}B{\to}C$, or $A{\to}B{\leftarrow}C$) embedded among 7 nuisance AR(1) channels ($\rho{=}0.8$). Coupling strengths $\alpha \sim \mathrm{Uniform}(0.3, 0.7)$; 1500 instances per class. The \texttt{pyspi} vocabulary yields $K{=}125$ statistics after variance filtering.

\textbf{Models.} Each baseline isolates a component: SPI-MPNN (full method); Correlation (fixed $|r_{ij}|$); SGC-Only (oracle-best single SPI); Fixed-SPI (fully connected, $\mathbf{E}_{ij}$ as edge features); Latent (learned node embeddings); Shuffled (SPIs permuted across pairs); Node-Only (no edge features).

% === FIGURE 1 ===
\begin{figure*}[t]
\vskip 0.2in
\begin{center}
\framebox[\textwidth]{\rule{0pt}{2.2in}\textbf{Figure 1 placeholder:} Sample efficiency curves --- F1 vs.\ $n$/class (log-scale $x$-axis), $\pm$1 s.d.\ bands.}
\caption{Sample efficiency on VAR(1) topology classification (10 seeds). SPI-vocabulary models achieve near-perfect accuracy while symmetric baselines plateau below $2/3$.}
\label{fig:sample_efficiency}
\end{center}
\vskip -0.2in
\end{figure*}

\textbf{Results.} Macro F1 (\%), 10 seeds. Top-$d{=}5$, group $\lambda{=}0.02$.

\begin{table}[t]
\caption{Main results (macro F1 \%, 10 seeds).}
\label{tab:main}
\vskip 0.1in
\begin{center}
\begin{small}
\begin{tabular}{lcccccc}
\toprule
$n$/cls & SPI & Fixed & SGC & Corr. & Lat. & Shuf. \\
\midrule
20   & 57{\scriptsize$\pm$22} & \textbf{72}{\scriptsize$\pm$12} & 61{\scriptsize$\pm$9}  & 35{\scriptsize$\pm$5}  & 30{\scriptsize$\pm$2} & 34{\scriptsize$\pm$4} \\
50   & 92{\scriptsize$\pm$2}  & 92{\scriptsize$\pm$2}  & \textbf{95}{\scriptsize$\pm$8}  & 47{\scriptsize$\pm$13} & 33{\scriptsize$\pm$1} & 38{\scriptsize$\pm$5} \\
100  & 97{\scriptsize$\pm$2}  & \textbf{98}{\scriptsize$\pm$1}  & 83{\scriptsize$\pm$17} & 57{\scriptsize$\pm$2}  & 32{\scriptsize$\pm$2} & 37{\scriptsize$\pm$3} \\
500  & 98{\scriptsize$\pm$4}  & \textbf{99}{\scriptsize$\pm$.5} & 87{\scriptsize$\pm$24} & 59{\scriptsize$\pm$1}  & 31{\scriptsize$\pm$2} & 55{\scriptsize$\pm$19}\\
1000 & \textbf{100}{\scriptsize$\pm$.2}& 99{\scriptsize$\pm$1} & 96{\scriptsize$\pm$11} & 59{\scriptsize$\pm$1}  & 31{\scriptsize$\pm$3} & 80{\scriptsize$\pm$13}\\
\bottomrule
\end{tabular}
\end{small}
\end{center}
\vskip -0.15in
\end{table}

Additional ablations (not shown): TE-only (29--18\%, chance), Top-3 oracle SPIs (67--100\%, std up to 27\%), MLP-Mix (63--99\%), Edge-Ablation (35--88\%).

\textbf{Interpretation.}

\textbf{(i) The SPI vocabulary is the dominant inductive bias --- and diversity provides stability.} All SPI-consuming models dramatically outperform non-SPI baselines. SGC-only (61--96\%, std 8--24\%), the oracle-best SPI, achieves high peak performance but with severe seed instability: 2/10 seeds catastrophically fail at $n{=}500$. Only the full vocabulary (std 0.2--2\% at $n \geq 50$, zero catastrophic failures) provides stable performance --- the 125 SPIs function as a multi-view representation whose diversity yields variance reduction.

\textbf{(ii) Symmetric construction is empirically limited, as predicted.} Correlation ($\leq$59\%) stays below the theoretical $2/3$ ceiling. The latent model (31\%) performs at chance, confirming that symmetric node-embedding construction cannot access the chain/fork distinction.

\textbf{(iii) Edge content is load-bearing.} The edge-ablation model achieves only 88$\pm$23\% at $n{=}500$ compared to 98$\pm$4\% for the full model. The MPNN actively uses the SPI descriptor $\phi(\mathbf{E}_{ij})$ to condition messages.

\textbf{(iv) The learned $\mathbf{w}$ is an interpretable statistical signature.} Under group lasso ($\lambda_g{=}0.02$), the causal family carries 6$\times$ the $L_2$ norm of the next family. The top-weighted SPIs are all spectral Granger causality and time-domain Granger causality measures, recovering the generating process.

%-----------------------------------------------------------------------
\section{Discussion}
\label{sec:discussion}
%-----------------------------------------------------------------------

Graph structure learning methods \citep{zhu2021} universally use latent embeddings. CauGNN \citep{duan2023} constructs graphs from a single pre-specified statistic; ADLGNN \citep{sriramulu2023} initialises from statistics but overwrites with attention, losing interpretability. Our approach occupies an unoccupied point: a \emph{vocabulary} of named statistics with \emph{learnable per-statistic weights} that remain inspectable after training. \citet{nguyen2025} argued that named statistical features capture interactions that raw-value methods miss --- we operationalise this insight within a GNN framework.

SPI computation is $O(K \cdot M^2)$ per sample; amortisable but limits scalability. All SPIs are bivariate; higher-order interactions require multi-hop message passing. The current study is synthetic; real-world validation (e.g., EEG motor imagery, fMRI) is needed.

The statistical signature produced by $\mathbf{w}$ --- identifying spectral Granger causality as the dominant coupling mode --- constitutes a testable hypothesis about the data-generating process. On real-world MTS, the same mechanism would identify which notion of coupling is relevant to the task, yielding domain-interpretable graph structure directly from the learned parameters.

\bibliography{references}
\bibliographystyle{icml2025}

\end{document}
